\chapter{Results and evaluation}
\label{ch:evaluation}

This chapter presents the measurement results and interprets them in terms of
the mechanisms described in the three preceding chapters. Unless stated
otherwise, the plots concern 32-bit integers and a uniform input
distribution. The effect of data type and distribution is examined separately
in sections~\ref{sec:eval-dtypes} and~\ref{sec:eval-distributions}.

\section{Correctness}
\label{sec:eval-correctness}

A precondition for any timing comparison is the correctness of the result.
The main sweep comprises 2700 cases, covering all three architectures, the
three algorithms, the five data types, input sizes from $2^3$ to $2^{24}$,
and the four values of k. All were verified against the ground truth of
section~\ref{sec:meth-correctness}, and all passed.

\begin{figure}[H]
    \centering
    \includegraphics[width=0.85\textwidth]{int/correctness/pass_rate}
    \caption{Pass rate of the correctness check for the two algorithms of this
    work, by architecture.}
    \label{fig:eval-pass-rate}
\end{figure}

This result is not a procedural detail. It means that every time reported in
what follows corresponds to an implementation that produced the correct set
of k elements, and that the comparisons that follow indeed concern the same
computation.

\section{Execution time}
\label{sec:eval-time}

\subsection{Dependence on input size}
\label{sec:eval-time-n}

\begin{figure}[H]
    \centering
    \begin{subfigure}{\textwidth}
        \centering
        \includegraphics[width=0.88\textwidth]{int/time/time_vs_n_algo_k8}
        \caption{Algorithmic scope.}
        \label{fig:eval-time-n-algo}
    \end{subfigure}

    \vspace{1em}

    \begin{subfigure}{\textwidth}
        \centering
        \includegraphics[width=0.88\textwidth]{int/time/time_vs_n_e2e_k8}
        \caption{End-to-end scope.}
        \label{fig:eval-time-n-e2e}
    \end{subfigure}
    \caption[Time as a function of size, under the two measurement scopes]{Time as a function of input size, for k equal to 8, under the two
    measurement scopes. The ranking of the architectures reverses completely
    between the two.}
    \label{fig:eval-time-n-k8}
\end{figure}

In the upper figure, which concerns the algorithmic scope, the GPU is the
fastest architecture across the entire range of sizes, from the smallest to
the largest. At the smallest input it completes in \qty{0.008}{\milli\second}
against \qty{0.048}{\milli\second} for the CPU, and at the largest in
\qty{0.42}{\milli\second} against \qty{1.41}{\milli\second}.

The other two architectures show a clear floor at small sizes. The NPU curves
remain completely flat around half a millisecond up to about $2^{16}$
elements: throughout this range the time does not depend on the input size at
all, since it is dominated by the fixed cost of dispatch and waiting. The CPU
shows an analogous but much lower floor, on the order of five hundredths of a
millisecond.

This picture, however, is completely reversed in the lower figure, once the
measurement includes data transfer as well. In the end-to-end scope the CPU is
faster than every other architecture at \emph{all} sizes examined, including
the largest: the GPU kernel's \qty{0.42}{\milli\second} becomes
\qty{5.81}{\milli\second} in total, while the CPU remains at
\qty{1.41}{\milli\second}, since for it the two scopes coincide.

This contrast is the first appearance of a pattern that recurs throughout the
chapter. The algorithmic scope measures how well each device performs the
computation; the end-to-end scope measures whether it is worth delegating to
it in the first place. The two answers do not coincide, and
section~\ref{sec:eval-energy} shows that it is the second that also
determines the energy picture.

The NPU requires \qty{10.46}{\milli\second} at the largest input, i.e. more
than any other architecture under either scope, for the reasons given in
section~\ref{sec:npu-bottlenecks}.

\begin{figure}[H]
    \centering
    \includegraphics[width=0.9\textwidth]{int/time/time_vs_n_backend_compare_bitonic_k8}
    \caption{The same algorithm across the three architectures, for k equal to 8.}
    \label{fig:eval-backend-bitonic}
\end{figure}

Comparing the same bitonic sort across the three architectures is
illuminating, because the same level schedule is executed and the difference
is due solely to the implementation. The GPU outperforms the CPU by a factor
of eight, as expected from the difference in available bandwidth. The NPU sits
close to the CPU, which is consistent with the finding that its time is spent
almost entirely on the host.

\subsection{Dependence on the parameter k}
\label{sec:eval-time-k}

The dependence on k is the point at which the ranking of the algorithms
reverses, and is perhaps the most significant finding of this work.

\begin{figure}[H]
    \centering
    \begin{subfigure}{\textwidth}
        \centering
        \includegraphics[width=0.88\textwidth]{int/time/time_vs_n_k_colored_cpu_bitonic}
        \caption{Bitonic sort, CPU.}
    \end{subfigure}

    \vspace{1em}

    \begin{subfigure}{\textwidth}
        \centering
        \includegraphics[width=0.88\textwidth]{int/time/time_vs_n_k_colored_gpu_map_reduce}
        \caption{Map-reduce, GPU.}
    \end{subfigure}
    \caption[Time as a function of size, by value of k]{Time as a function of input size, with a separate curve for each
    value of k. The bitonic sort curves remain close to one another, while the
    map-reduce curves diverge by orders of magnitude.}
    \label{fig:eval-k-colored}
\end{figure}

The juxtaposition in figure~\ref{fig:eval-k-colored} captures the phenomenon
directly. For bitonic sort the curves of the four values of k nearly
coincide, since k affects only the size of the sorted blocks. For map-reduce
on the GPU the same curves are separated by three orders of magnitude.

\begin{table}[H]
    \centering
    \begin{tabular}{lrrrr}
        \hline
        & \multicolumn{2}{c}{k=8} & \multicolumn{2}{c}{k=131072} \\
        Implementation & ms & rank & ms & rank \\
        \hline
        GPU map-reduce & 0.42 & 1 & 1191.44 & 6 \\
        GPU bitonic & 1.22 & 2 & 8.00 & 1 \\
        CPU map-reduce & 1.41 & 3 & 62.45 & 5 \\
        NPU bitonic & 8.79 & 4 & 56.99 & 3 \\
        CPU bitonic & 10.10 & 5 & 47.67 & 2 \\
        NPU map-reduce & 10.46 & 6 & 60.30 & 4 \\
        \hline
    \end{tabular}
    \caption[Algorithmic time for the two extremes of k]{Algorithmic time for the two extremes of k, at an input of $2^{24}$
    elements. The ranking concerns the six implementations of this work and
    does not include the ground-truth implementations.}
    \label{tab:eval-k-inversion}
\end{table}

The fastest implementation for small k becomes the slowest for large k, a
difference of three orders of magnitude. GPU map-reduce slows down by a
factor of 2800 between the two extremes, whereas bitonic sort on the same
architecture slows down by barely a factor of seven.

The reversal is not a peculiarity of this implementation. Independent studies
on GPUs from two different vendors report the same dependence on k and
address it with adaptive algorithm switching~\cite{zhang2023gputopk,lin2026adaptivetopk};
there, the growth in time of partial-sorting methods with k is attributed
explicitly to the $O(\log^2 k)$ complexity of the underlying bitonic network.

The cause was given in section~\ref{sec:gpu-mapreduce-limit}: the heap only
filters once it is full, so the number of threads that contribute is limited
to the ratio $n/k$. For the sizes in table~\ref{tab:eval-k-inversion} this
ratio is 128, i.e. the device effectively employs one hundred and twenty-eight
threads out of the thousands it has available. The CPU, with only eight
threads, never runs into this same limit, and for that reason its slowdown is
markedly milder.

\begin{figure}[H]
    \centering
    \includegraphics[width=0.9\textwidth]{int/time/time_vs_n_algo_compare_gpu_k131072}
    \caption{The three algorithms on the GPU for k equal to 131072.}
    \label{fig:eval-algo-compare-gpu}
\end{figure}

The practical conclusion is that there is no universally preferable
algorithm. The choice between a sorting network and heap-based selection
depends on the ratio $n/k$ and not on the architecture.

\section{Comparison across architectures}
\label{sec:eval-backends}

\begin{figure}[H]
    \centering
    \begin{subfigure}{0.49\textwidth}
        \centering
        \includegraphics[width=\textwidth]{int/time/speedup_vs_gt_cpu_k8}
        \caption{CPU, k=8.}
    \end{subfigure}
    \hfill
    \begin{subfigure}{0.49\textwidth}
        \centering
        \includegraphics[width=\textwidth]{int/time/speedup_vs_gt_cpu_k131072}
        \caption{CPU, k=131072.}
    \end{subfigure}

    \vspace{0.5em}

    \begin{subfigure}{0.49\textwidth}
        \centering
        \includegraphics[width=\textwidth]{int/time/speedup_vs_gt_gpu_k8}
        \caption{GPU, k=8.}
    \end{subfigure}
    \hfill
    \begin{subfigure}{0.49\textwidth}
        \centering
        \includegraphics[width=\textwidth]{int/time/speedup_vs_gt_gpu_k131072}
        \caption{GPU, k=131072.}
    \end{subfigure}

    \vspace{0.5em}

    \begin{subfigure}{0.49\textwidth}
        \centering
        \includegraphics[width=\textwidth]{int/time/speedup_vs_gt_npu_k8}
        \caption{NPU, k=8.}
    \end{subfigure}
    \hfill
    \begin{subfigure}{0.49\textwidth}
        \centering
        \includegraphics[width=\textwidth]{int/time/speedup_vs_gt_npu_k131072}
        \caption{NPU, k=131072.}
    \end{subfigure}
    \caption[Speedup relative to the ground-truth implementation]{Speedup relative to the ground-truth implementation of the same
    architecture. Left column k equal to 8, right column k equal to 131072,
    one row per architecture. Values above unity mean that this work's
    implementation is faster than the library.}
    \label{fig:eval-speedup}
\end{figure}

Comparing each implementation with the reference of the same architecture
gives the most honest picture, since it neutralises hardware differences.

On the CPU, map-reduce is faster than the library by a factor of three for
small k. It should be recalled, however, that the limitation of
section~\ref{sec:cpu-reference} applies: the library is single-threaded while
this implementation uses eight threads, so this factor is not purely
algorithmic.

The largest speedup in the entire figure appears elsewhere, however, and is a
consequence of the reversal in section~\ref{sec:eval-time-k}. For large k,
where heap-based selection collapses, bitonic sort on the CPU reaches nearly
seven times faster than the library. The reason is that the library, too,
relies on a heap of size k and suffers the same overhead, while the sorting
network remains essentially unaffected. The advantage weakens as the input
grows further and the network's data traffic starts to dominate, narrowing to
just above unity at the largest input.

On the GPU the comparison is strict, since the library runs on the same
device with the same parallelism. There, map-reduce marginally outperforms
for k equal to 8, while for large k the library is faster by more than a
thousand times. Radix selection has no point of collapse with respect to k,
as explained in section~\ref{sec:gpu-reference-cub}.

On the NPU, for the integer types to which the figure refers, no
implementation surpasses the library in any parameter region. This result is
expected rather than disappointing: as shown in section~\ref{sec:npu-bottlenecks},
the very processor that would run the library is in any case consumed by
preparing the data for the device.

The qualifier ``for the integer types'' is not a hedge. For floating-point
types, and especially for half precision, the picture reverses and the NPU
does indeed surpass the library; the explanation is given in
section~\ref{sec:eval-dtypes}, after the behaviour across types is first
presented.

\section{Roofline analysis}
\label{sec:eval-roofline}

\begin{figure}[H]
    \centering
    \begin{subfigure}{\textwidth}
        \centering
        \includegraphics[width=0.86\textwidth]{int/memory/roofline_kernels_k8}
        \caption{k equal to 8.}
        \label{fig:eval-roofline-k8}
    \end{subfigure}

    \vspace{1em}

    \begin{subfigure}{\textwidth}
        \centering
        \includegraphics[width=0.86\textwidth]{int/memory/roofline_kernels_k131072}
        \caption{k equal to 131072.}
        \label{fig:eval-roofline-k131072}
    \end{subfigure}
    \caption[Placement of the runs against the measured ceilings]{Placement of the runs against the measured ceilings, for the largest
    input. Restricted to the maximum size, where the working set exceeds the
    cache and the main-memory ceiling is the relevant one.}
    \label{fig:eval-roofline}
\end{figure}

The diagram places each run at its level of arithmetic intensity and
achieved throughput. All implementations cluster in the left-hand region,
i.e. at very low arithmetic intensity, and none approaches the horizontal
part of the ceiling. No implementation of this work, on any architecture, is
limited by compute capability. This finding agrees with an independent
measurement on data-center accelerators, where the fastest known selection
kernel utilises over ninety percent of the available bandwidth against about
half of compute capability, and where time scales with each device's
bandwidth~\cite{zhang2023gputopk}.

Juxtaposing the two values of k adds two observations. In both cases bitonic
sort on the GPU sits practically on the sloped section of its own ceiling, which
visually confirms the finding of section~\ref{sec:gpu-bottlenecks}: the
implementation utilises the available bandwidth almost fully and is limited
exclusively by the volume it is obliged to move.

The second observation concerns map-reduce on the same architecture. For k
equal to 8 it sits near the other points, while for the large value it
collapses by three orders of magnitude and moves away from both ceilings
equally. This position, far from both the sloped and the horizontal section,
is the signature of a limitation that is neither memory nor compute, and
corresponds to the collapse of parallelism described in
section~\ref{sec:gpu-mapreduce-limit}.

This finding is structural rather than incidental. Top-k selection performs
essentially one comparison per element and reuses almost nothing, so its
arithmetic intensity is low by definition.

\begin{figure}[H]
    \centering
    \includegraphics[width=0.85\textwidth]{machine/cache_ladder}
    \caption{Measured read bandwidth as a function of working-set size, for
    the three architectures.}
    \label{fig:eval-cache-ladder}
\end{figure}

The ceiling against which each run is compared is not fixed but a function of
its working set, per section~\ref{sec:meth-roofline}. Figure~\ref{fig:eval-cache-ladder}
shows the measured ladder: the steps correspond to the levels of the memory
hierarchy, and the drop toward the right-hand edge marks the point where the
working set stops fitting in the last level.

\subsection{Memory bandwidth utilisation}
\label{sec:eval-bandwidth}

\begin{figure}[H]
    \centering
    \includegraphics[width=0.9\textwidth]{int/memory/cpu_roof_utilization_k8}
    \caption{Utilisation of the two ceilings on the CPU, for k equal to 8.}
    \label{fig:eval-roof-util}
\end{figure}

\begin{figure}[H]
    \centering
    \begin{subfigure}{\textwidth}
        \centering
        \includegraphics[width=0.66\textwidth]{int/memory/cpu_memory_bandwidth_vs_n_k8}
        \caption{CPU.}
    \end{subfigure}

    \vspace{0.6em}

    \begin{subfigure}{\textwidth}
        \centering
        \includegraphics[width=0.66\textwidth]{int/memory/gpu_memory_bandwidth_vs_n_k8}
        \caption{GPU.}
    \end{subfigure}

    \vspace{0.6em}

    \begin{subfigure}{\textwidth}
        \centering
        \includegraphics[width=0.66\textwidth]{int/memory/npu_memory_bandwidth_vs_n_k8}
        \caption{NPU.}
    \end{subfigure}
    \caption[Achieved bandwidth by architecture]{Achieved bandwidth as a function of input size, one architecture per
    panel, with the measured ceiling drawn as a dashed line. The distance of
    each curve from the dashed line is the unexploited margin.}
    \label{fig:eval-bandwidth}
\end{figure}

Utilisation of the memory ceiling increases with input size and stabilises once
the working set exceeds the cache. At the largest input, bitonic sort on the
CPU reaches two-thirds of the ceiling, and on the GPU a comparable percentage,
while utilisation of the compute ceiling remains below five percent.

The juxtaposition of bitonic sort with the library on the GPU deserves
particular attention, because it shows that proximity to the ceiling does not
equate to speed. Bitonic sort clearly utilises a larger fraction of the
available bandwidth than the library, and is nonetheless slower, because it
moves seven times the volume of data. Utilisation measures the quality of the
implementation; total volume measures the quality of the algorithm.

\section{Energy assessment}
\label{sec:eval-energy}

\begin{figure}[H]
    \centering
    \includegraphics[width=0.9\textwidth]{int/time/time_vs_n_metric_compare_gpu}
    \caption[Algorithmic and end-to-end time on the GPU]{Algorithmic and end-to-end time on the GPU. The distance between the
    two curves of each algorithm is the cost of transfers and kernel launches.}
    \label{fig:eval-metric-compare}
\end{figure}

Before presenting the energy measurements, figure~\ref{fig:eval-metric-compare}
is useful, since it shows why the two measurement scopes diverge so sharply
on the GPU. The end-to-end curve remains steady near three tenths of a
millisecond up to about $2^{16}$ elements: throughout this range the kernel
is essentially invisible and the time is spent entirely on transfers and
launches. Even at the largest input, the end-to-end time remains roughly
three times the algorithmic time.

\begin{figure}[H]
    \centering
    \includegraphics[width=0.9\textwidth]{int/energy/energy_by_backend_net_k8}
    \caption{Net energy by architecture and algorithm, for k equal to 8.}
    \label{fig:eval-energy}
\end{figure}

The energy picture differs substantially from the time picture, and the
difference is due to the symmetric measurement scope of
section~\ref{sec:meth-energy-scope}: each architecture is charged the total
cost of solving the problem, including the host.

At the largest input and for small k, CPU map-reduce consumes about
\qty{20}{\milli\joule} per run and is the most economical implementation. The
corresponding implementations on the GPU consume from \qty{129}{\milli\joule}
to \qty{150}{\milli\joule}, i.e. up to seven times more, despite being faster
in the algorithmic scope. The most expensive implementation of all is bitonic
sort on the GPU, while the corresponding CPU implementation costs
\qty{141}{\milli\joule}.

The apparent contradiction is explained by section~\ref{sec:gpu-memory}. The
GPU's algorithmic scope includes only the kernel, while the end-to-end scope
includes the transfers, which consume time and power on both sides of the
bus. A device that is four times faster at the kernel but needs twice the
kernel's own time just to receive the data, while simultaneously drawing
board power, is not more economical.

\begin{figure}[H]
    \centering
    \includegraphics[width=0.9\textwidth]{int/energy/power_by_backend_net_k8}
    \caption{Net power by architecture and algorithm, for k equal to 8.}
    \label{fig:eval-power}
\end{figure}

Juxtaposing figures~\ref{fig:eval-energy} and~\ref{fig:eval-power} highlights
a distinction that is easy to conflate, and one that particularly concerns
the NPU. This device has the \emph{lowest} power draw of all
implementations, with map-reduce measuring \qty{8.1}{\watt} against
\qty{8.5}{\watt} for the corresponding CPU implementation and
\qty{18.6}{\watt} for the GPU.

Nonetheless it consumes five times the energy of CPU map-reduce,
\qty{0.097}{\joule} against \qty{0.020}{\joule}. The reason is that energy is
the product of power and time, and the NPU needs roughly seven times the time
for the same task. The five-percent advantage in power is wiped out by the
sevenfold disadvantage in duration.

This observation generalises: low power does not imply low energy. The
principle that moving data costs more energy than the computation itself
does~\cite{chen2016eyeriss} holds here as well; what is missing is the reuse
that would amortise it. A low-power architecture pays off energetically only
if it completes the task in comparable time, and for this particular problem
that is not the case.

\begin{figure}[H]
    \centering
    \includegraphics[width=0.9\textwidth]{int/energy/time_vs_energy}
    \caption[End-to-end time and energy per measured case]{End-to-end time and total energy of each measured case. Both
    quantities are minimised, so the preferable cases lie toward the
    bottom-left corner.}
    \label{fig:eval-time-energy}
\end{figure}

Figure~\ref{fig:eval-time-energy} places each measured case in a plane of two
criteria that are minimised simultaneously. A case is said to be dominated by
another when the latter is no worse in either quantity and is better in at
least one; such a case is never worth choosing. The set of cases dominated by
none other constitutes the Pareto front, i.e. exactly those choices in which
improving one criterion requires giving up ground on the other.

The result is illuminating precisely because it is negative. The
measurements form a narrow diagonal band, with a correlation coefficient
between the logarithms of time and energy equal to 0.97. Of the 315 cases,
only one belongs to the front. That is, no substantial trade-off exists
between speed and energy for this problem: the fastest configuration is also
the most economical.

The reason follows from the two preceding figures. Power varies among
implementations by barely a factor of 2.3, while time varies by four orders
of magnitude. When one factor of the product is nearly constant, the product
follows the other. The possibility of trading longer duration for lower
power, which on other systems is a genuine design lever, does not exist here.

Finally, it is worth noting which cases make up the front: all of them run on
the CPU. No configuration of the GPU or the NPU is Pareto-optimal at any
point of the parameter space examined.

\subsection{Energy-delay product}
\label{sec:eval-edp}

\begin{figure}[H]
    \centering
    \includegraphics[width=0.9\textwidth]{int/energy/edp_vs_n_net_k8}
    \caption{Energy-delay product as a function of input size.}
    \label{fig:eval-edp}
\end{figure}

The energy-delay product condenses the two criteria into one and rewards
implementations that are simultaneously fast and economical. The resulting
ranking is the clearest in the entire chapter. CPU map-reduce has the lowest
product of any implementation, with the next-lowest, GPU map-reduce, trailing
by a factor of twenty-two, and the worst by a factor of eighty-eight.

The gap widens further as the input shrinks: at the medium sizes of the
figure, no other implementation comes within less than two orders of
magnitude.

This result is the most counterintuitive of those presented here. For top-k
selection with small k, offloading to an accelerator does not merely fail to
pay off energetically; it is clearly worse than not offloading at all.

\section{Effect of data type}
\label{sec:eval-dtypes}

\begin{table}[H]
    \centering
    \begin{tabular}{lrrrrr}
        \hline
        Implementation & half & int & uint & float & double \\
        \hline
        CPU bitonic & 5.67 & 10.10 & 10.57 & 10.18 & 17.89 \\
        CPU map-reduce & 0.48 & 1.41 & 1.39 & 1.53 & 2.47 \\
        CPU reference & 13.34 & 4.71 & 4.73 & 5.61 & 5.63 \\
        GPU bitonic & 0.75 & 1.22 & 1.25 & 1.20 & 2.73 \\
        NPU bitonic & 8.18 & 8.79 & 8.91 & 8.77 & 28.36 \\
        \hline
    \end{tabular}
    \caption[Algorithmic time by data type]{Algorithmic time by data type, in ms, for $2^{24}$ elements and k
    equal to 8.}
    \label{tab:eval-dtypes}
\end{table}

This work's implementations track the width of the type, as expected of
algorithms limited by memory. Half precision is consistently the fastest,
with CPU bitonic sort completing in just over half the time of single
precision, precisely because it moves half the bytes at the same rate of
comparisons.

The reference implementation behaves oppositely: half precision is its
slowest case. The difference reveals that the library is not limited by
memory but by the cost of the scalar comparison, which for half precision
includes a conversion. The ranking among types is therefore not a property of
the hardware but of the implementation.

Double precision on the NPU stands out, with three times the time of any
other type. The overhead is due entirely to the additional refinement pass of
section~\ref{sec:npu-refinement}, which is required because the 32-bit key is
not sufficient to express 64-bit values.

The combination of these two opposite behaviours produces the single point in
the entire work where the NPU outperforms its own library. For half
precision, bitonic sort on the NPU is faster than the library in twenty-eight
of the measured cases, with the largest gap reaching a factor of two and a
half. The phenomenon appears more weakly in the other two floating-point
types, and never appears for the integer types.

The cause is not that the NPU handles half precision well, but that the
library handles it poorly. The array's kernel operates on 32-bit integers
regardless of type, so the cost of the type shifts to host encoding, which
for half precision is vectorised per section~\ref{sec:npu-encoding}. The
library, by contrast, performs scalar comparisons with an intermediate
conversion. Offloading thus proves advantageous precisely where the
established solution is weakest, not where the NPU is strongest.

\section{Sensitivity to the input distribution}
\label{sec:eval-distributions}

\begin{table}[H]
    \centering
    \begin{tabular}{lrrrr}
        \hline
        Implementation & uniform & adversarial & sorted & reverse \\
        \hline
        CPU bitonic & 44.50 & 44.75 & 44.57 & 44.41 \\
        GPU bitonic & 7.99 & 8.07 & 7.96 & 7.96 \\
        CPU reference & 74.03 & 73.91 & 766.66 & 10.65 \\
        NPU bitonic & 57.22 & 56.95 & 679.00 & 13.20 \\
        GPU map-reduce & 1192.75 & 1193.12 & 1141.87 & 1109.16 \\
        \hline
    \end{tabular}
    \caption[Algorithmic time by input distribution]{Algorithmic time by input distribution, in ms, for single precision,
    $2^{24}$ elements, and k equal to 131072.}
    \label{tab:eval-dists}
\end{table}

\begin{figure}[H]
    \centering
    \begin{subfigure}{\textwidth}
        \centering
        \includegraphics[width=\textwidth]{float/dist/dist_sensitivity_k8}
        \caption{k=8.}
        \label{fig:eval-dists-k8}
    \end{subfigure}

    \vspace{1em}

    \begin{subfigure}{\textwidth}
        \centering
        \includegraphics[width=\textwidth]{float/dist/dist_sensitivity_k131072}
        \caption{k=131072.}
        \label{fig:eval-dists-k131072}
    \end{subfigure}
    \caption[Sensitivity of time to the input distribution]{Sensitivity of time to the input distribution, for single precision
    and $2^{24}$ elements. Each bar gives the time of one distribution
    relative to the time of the uniform distribution for the same
    implementation; unity means complete independence from the data.}
    \label{fig:eval-dists}
\end{figure}

Normalising against the uniform distribution allows all nine combinations to
be compared on the same figure, even though their absolute times differ by
orders of magnitude. An implementation unaffected by the data appears as five
equal bars of height one, while a sensitive one appears with bars diverging
in both directions.

The figure first confirms the expected: bitonic sort on the CPU and the GPU
is practically unaffected by the distribution at both values of k, with
overall variation below a factor of 1.12. This property makes it a reliable
witness: if its time were to change, the change would be attributable to the
measurement rather than to the algorithm.

Heap-based implementations show the opposite extreme, and indeed regardless
of who wrote the code. The reference library varies by a factor of
seventy-two between ascending and descending input, but so does our own CPU
map-reduce, by a factor of twenty-five, through the same mechanism described
in section~\ref{sec:cpu-mapreduce}. Sensitivity to the data is therefore a
property of the algorithm rather than of the implementation: whoever chooses
a heap inherits it.

Two observations deserve particular attention. First, bitonic sort on the NPU
\emph{is} sensitive to the distribution, unlike the other two architectures,
with a variation factor of fifty-one. This sensitivity does not originate in
the array, which executes a fixed network, but in the heap-based merge
performed by the host, which, as shown in section~\ref{sec:npu-pipeline},
dominates the time.

Second, juxtaposing the two values of k isolates the cause of the
sensitivity. GPU map-reduce is sensitive for k equal to 8, where the sorted
input slows it down by a factor of 4.2, but becomes nearly unaffected for k
equal to 131072, where the same input changes it by only 4\%. The same code,
with the same data, changes behaviour because the limiting factor changes: in
the second case it is not rejection that limits it but the lack of
parallelism described in section~\ref{sec:gpu-mapreduce-limit}. When the
limiting factor is not the data, their distribution ceases to matter.

The measurements presented cover the entire parameter space of
section~\ref{sec:meth-parameters} and agree with one another as to which
factor limits each implementation. Their interpretation, the limits within
which they hold, and the directions they open up are examined in
chapter~\ref{ch:conclusions}.
